% ===========================================================================
%  main.tex — 厦门大学课程论文 主文件
%  Xiamen University Course Paper — Main Document
% ===========================================================================
\documentclass{xmu_course}

% ===========================================================================
%  1. 封面信息  Cover Information
%     以下命令用于 \makecover 生成封面
% ===========================================================================
\covertitle{2023--2024学年\ 春季学期期末论文}  % 封面大标题
\coursename{自然辩证法概论}                     % 课程名称
\classinfo{周一\ 第3--4节}                      % 班级 / 上课时间
\college{人文学院}                              % 所在学院
\major{哲学}                                    % 专业
\grade{2023级}                                  % 年级
\authorname{张\ 三}                             % 姓名
\studentid{12320230001}                         % 学号
\coverdate{2024年6月}                           % 封面日期

% ===========================================================================
%  2. 正文元信息  Title Block Data
%     以下命令用于 \maketitle 输出篇首页
%     （若不需要单独篇首页，可全部注释，正文直接从 section 开始）
% ===========================================================================
\title{论文标题}                                % 中文篇名（二号宋体，≤20字）
\author{张\ 三}                                % 作者姓名（四号楷体）
\affiliation{厦门大学\ 人文学院，福建\ 厦门\ 361005}  % 作者单位（小五号宋体）

\clcnumber{B82-057}                             % 中图分类号
\doccode{A}                                     % 文献标识码

% 中文摘要（200--500字，小五号宋体）
\abstract{%
在这里输入中文摘要，一般为200--500字。摘要应能反映论文的主要信息，
不应出现"本文""作者认为"等字样。%
}

% 中文关键词（3--5个，用分号隔开，小五号宋体）
\keywords{关键词1；关键词2；关键词3；关键词4}

% 收稿日期 / 基金项目 / 作者简介（首页地脚，小五号楷体）
\receivedate{2024-05-20}
\fundproject{基金项目名称（项目编号）}
\authorbio{张三，男，籍贯，厦门大学人文学院教授，哲学博士}

% ===========================================================================
%  3. 英文信息（放在正文末尾，参考文献之后）
% ===========================================================================
\entitle{English Title Here}
\enauthor{ZHANG San}
\enaffiliation{(School of Humanities, Xiamen University, Xiamen 361005, Fujian, China)}
\enabstract{English abstract text goes here. It should correspond to the
Chinese abstract and be free of pragmatic errors.}
\enkeywords{keyword1; keyword2; keyword3; keyword4}

% ===========================================================================
%  文档开始
% ===========================================================================
\begin{document}

% -------- 封面 --------
\makecover

% -------- 篇首页（标题、作者、单位、摘要、关键词、地脚信息）--------
%           若不需要单独篇首页，注释掉下面两行即可
\maketitle
\makefirstpagefooter

% ===========================================================================
%  正文
% ===========================================================================
\section{一级标题示例}

正文内容从这里开始。正文采用五号宋体字，单倍行距。版面为42字乘以42行。

\subsection{二级标题示例}

二级标题采用五号黑体字。更多的正文内容\footnote{这是脚注示例，脚注采用小五号宋体字，
列在当页地脚，使用阿拉伯数字圈码编序。}。

\section{另一个一级标题}

\subsection{又一个小节}

\begin{table}[htbp]
  \centering
  \caption{三线表示例}
  \label{tab:example}
  \begin{tabular}{@{}ccc@{}}
    \tabletoprule
    列一 & 列二 & 列三 \\
    \tablemidrule
    A & B & C \\
    D & E & F \\
    \tablebottomrule
  \end{tabular}
\end{table}

% ===========================================================================
%  参考文献
% ===========================================================================
\begin{thebibliography}{99}
\bibitem{ref1}
作者. 文献标题[J]. 期刊名, 年份, 卷(期): 起止页码.

\bibitem{ref2}
作者. 文献标题[M]. 出版地: 出版社, 年份.
\end{thebibliography}

% ===========================================================================
%  英文摘要（放在文后）
% ===========================================================================
\makeenabstract

\end{document}
